% ---
% id: transmitter/lem:ensemble-presentation
% kind: lem
% title: "Certificate defined by an ensemble presentation"
% summary: >-
%   A virtual ensemble in sector blocks defines, by canonical purification and
%   a square-root measurement, a block certificate with pure blocks that
%   exactly reproduces Alice's marginal and, under that measurement, the
%   ensemble. Only her marginal, a transposed sum of declared blocks, is formed
%   as numbers.
% prominence: main
% status: implemented
% uses: [def:srcert, def:source-spec]
% mentions: [thm:certificate-sufficiency]
% assumes: []
% refs: []
% provenance: newly-derived
% implements: certiqs_key.transmitter.presentation:alice_marginal_blocks
% lean:
% notes: "The square roots in the statement are real-number objects of the proof and are formed by no consumer. Part (v) is what a consumer's ensemble-form declaration reads; that form is the consumer's to define and is not defined here."
% ---

The ensemble presentation of Definition~\ref{def:srcert} declares only the
weighted virtual ensemble, in blocks along a sector structure of the virtual
transmitted system. The lemma supplies the objects it defines and shows that
they form a block-form certificate, so that no square root or quotient of a
declared number ever has to be declared; of the defined objects only Alice's
marginal is formed as numbers, that is, computed entry by entry.

\begin{lemma}\label{lem:ensemble-presentation}
Let the virtual transmitted system carry a sector structure
$\Hs{\Avirt}\supseteq\bigoplus_{t\in\mathbb{T}}\Hs{S_t}$, each $S_t$ spanned
by vectors $(\ket{i}_{S_t})_i$ of the declared basis, and let
$(\sigma^{\mathrm v,(t)}_x)_{t\in\mathbb{T}}$, $x\in\Xset$, be block lists
with $\sigma^{\mathrm v,(t)}_x\in\Pos{S_t}$ and
$\sum_{x}\sum_t\tr\sigma^{\mathrm v,(t)}_x=1$. For every $t$ let $A_t$ be a
copy of $S_t$ with basis $(\ket{i}_{A_t})_i$ in the same order, let
transposition $X\mapsto X^{\transpose}$ be taken in these bases, and put
\begin{equation}\label{eq:enspres}
\begin{gathered}
  \tau^{(t)}_{S}\coloneqq\sum_{x\in\Xset}\sigma^{\mathrm v,(t)}_x ,\qquad
  \tau^{(t)}_A\coloneqq\big(\tau^{(t)}_{S}\big)^{\transpose} ,\\
  \ket{\psi_t}\coloneqq\sum_i\ket{i}_{A_t}\otimes
     \big(\tau^{(t)}_{S}\big)^{1/2}\ket{i}_{S_t} ,\qquad
  \Psi_t\coloneqq\dyad{\psi_t} ,\\
  M^{(t)}_x\coloneqq\big(\tau^{(t)}_A\big)^{-1/2}
     \big(\sigma^{\mathrm v,(t)}_x\big)^{\transpose}
     \big(\tau^{(t)}_A\big)^{-1/2}
     +\delta_{x,x_0}\big(\idon{A_t}-P_t\big) ,
\end{gathered}
\end{equation}
where $P_t$ is the projector onto the support of $\tau^{(t)}_A$, the inverse
square root is taken on that support and is zero off it, and $x_0$ is the
first letter of $\Xset$ in its declared order. Then, with
$A\coloneqq\bigoplus_tA_t$:
\begin{enumerate}[label=(\roman*),leftmargin=2.2em]
\item \emph{(a certificate in block form)} every $\Psi_t$ is positive
semidefinite, $\sum_t\tr\Psi_t=1$, $\{M_x\}_x$ with
$M_x=\bigoplus_tM^{(t)}_x$ is a POVM presented as a block list,
$\tau_A=\bigoplus_t\tau^{(t)}_A$ is a state presented as a block list, and
both identities of \eqref{eq:srblock} hold: $(\Psi_t)_t$ satisfies every
condition Definition~\ref{def:srcert} places on a block-form certificate
for $(A,\mathbb{T},\tau_A,\{M_x\})$ against the virtual ensemble with
weighted members $\sigma^{\mathrm v}_x=\bigoplus_t\sigma^{\mathrm v,(t)}_x$,
(SR3) apart, which is an attestation;
\item \emph{(pure blocks)} every non-zero $\Psi_t$ is
$g_t\coloneqq\tr\tau^{(t)}_S$ times a rank-one projector, so the converse of
Theorem~\ref{thm:certificate-sufficiency}\,(ii) holds of this certificate;
\item \emph{(what is formed as numbers)} the entry $(i,j)$ of
$\tau^{(t)}_A$ is the entry $(j,i)$ of the sum $\tau^{(t)}_S$ of the
declared blocks. No
other object of \eqref{eq:enspres} enters (i) or (ii) through its entries:
$\Psi_t$ and $M^{(t)}_x$ contain square roots and quotients of the declared
numbers and are, in general, not declarable;
\item \emph{(the completion is immaterial)} the letter $x_0$ enters the
POVM condition and nothing else: (SR1), (SR2), (ii) and (v) hold with the
term $\idon{A_t}-P_t$ assigned to any letter, or split among the letters by
any positive semidefinite operators on the kernel of $\tau^{(t)}_A$ that
sum to $\idon{A_t}-P_t$;
\item \emph{(Alice's side without a square root)}
$\big(\tau^{(t)}_A\big)^{1/2}M^{(t)}_x\big(\tau^{(t)}_A\big)^{1/2}
=\big(\sigma^{\mathrm v,(t)}_x\big)^{\transpose}$ for every $x$ and $t$: the
operators $N^{(t)}_x\coloneqq(\sigma^{\mathrm v,(t)}_x)^{\transpose}$ are
entrywise the declared numbers, and
$\sum_xN^{(t)}_x=\tau^{(t)}_A$.
\end{enumerate}
\end{lemma}

\begin{proof}
\emph{Idea.} Each block is the canonical purification of the summed
ensemble of its sector, written through the maximally entangled vector; three
facts about transposition reduce every claim to an identity between the
declared blocks and their sum.

\emph{Preliminaries.} Fix $t$ and drop it from the notation; write
$K\coloneqq\tau_S^{1/2}$,
$\ket{\Omega}\coloneqq\sum_i\ket{i}_A\ket{i}_S$, so that
$\ket{\psi}=(\idon{A}\otimes K)\ket{\Omega}$, and let $Q$ be the projector
onto the support of $\tau_S$. Three facts are used. (a)
$\tr_A\big[(Y\otimes\idon{S})\dyad{\Omega}\big]=Y^{\transpose}$ for
$Y\in\Lin{A}$: with $\dyad{\Omega}=\sum_{ij}\ket{i}\!\bra{j}\otimes
\ket{i}\!\bra{j}$ the left-hand side is $\sum_{ij}\bra{j}Y\ket{i}\,
\ket{i}\!\bra{j}$, whose $(i,j)$ entry is $Y_{ji}$. (b)
$(\idon{}\otimes Z)\ket{\Omega}=(Z^{\transpose}\otimes\idon{})\ket{\Omega}$,
by comparing coefficients. (c) Transposition is linear, reverses products,
commutes with the adjoint and maps positive semidefinite operators to
positive semidefinite operators with the same spectrum; hence
$f(X^{\transpose})=f(X)^{\transpose}$ for every real function $f$ on the
spectrum of a Hermitian $X$. In particular $(\tau_A)^{1/2}=K^{\transpose}$,
$(\tau_A)^{-1/2}=(\tau_S^{-1/2})^{\transpose}$ and $P=Q^{\transpose}$.

\emph{(SR1).} By (b),
$\Psi=(K^{\transpose}\otimes\idon{})\dyad{\Omega}
(K^{\transpose}\otimes\idon{})^{\dagger}$, and
$\tr_S\dyad{\Omega}=\idon{A}$, so
$\tr_S\Psi=K^{\transpose}(K^{\transpose})^{\dagger}=(K^{\transpose})^2
=(K^2)^{\transpose}=\tau_S^{\transpose}=\tau_A$, the operator
$K^{\transpose}$ being Hermitian by (c).

\emph{(SR2).} $\tr_A\big[(M_x\otimes\idon{})\Psi\big]
=K\,\tr_A\big[(M_x\otimes\idon{})\dyad{\Omega}\big]K=K\,M_x^{\transpose}K$
by (a), $K$ acting on $S$ alone. By (c),
$M_x^{\transpose}=\tau_S^{-1/2}\sigma^{\mathrm v}_x\tau_S^{-1/2}
+\delta_{x,x_0}(\idon{S}-Q)$, and $K\tau_S^{-1/2}=Q$, $K(\idon{}-Q)=0$, so
$K\,M_x^{\transpose}K=Q\,\sigma^{\mathrm v}_xQ$. Finally
$0\preceq\sigma^{\mathrm v}_x\preceq\tau_S$ puts the support of
$\sigma^{\mathrm v}_x$ inside that of $\tau_S$, so
$Q\,\sigma^{\mathrm v}_xQ=\sigma^{\mathrm v}_x$.

\emph{The remaining conditions of (i).} $\Psi=\dyad{\psi}\succeq0$ and
$\tr\Psi=\tr\tr_S\Psi=\tr\tau_A=\tr\tau_S$, whose sum over $t$ is
$\sum_x\sum_t\tr\sigma^{\mathrm v,(t)}_x=1$. Each $M_x$ is positive
semidefinite: $(\sigma^{\mathrm v}_x)^{\transpose}\succeq0$ by (c), a
conjugation by the Hermitian operator $\tau_A^{-1/2}$ preserves positivity,
and $\idon{}-P\succeq0$. Their sum is
$\tau_A^{-1/2}\,\tau_S^{\transpose}\,\tau_A^{-1/2}+(\idon{}-P)
=P+(\idon{}-P)=\idon{A}$, by linearity of transposition. $\tau_A$ is
positive semidefinite by (c), and $\sum_t\tr\tau^{(t)}_A=1$ was just shown.
The weighted members $\bigoplus_t\sigma^{\mathrm v,(t)}_x$ of the virtual
ensemble are block diagonal along $(S_t)_t$ by their presentation as block
lists, with the declared blocks as their diagonal blocks, which is the
condition the block form places on the virtual ensemble; together with
\eqref{eq:srblock}, shown above sector by sector, this is (SR1) and (SR2) of
$\Psi=\bigoplus_t\Psi_t$ by Definition~\ref{def:srcert}.

(ii) $\Psi=\dyad{\psi}$ with $\braket{\psi}{\psi}=\tr\Psi=g$; where $g>0$,
$\Psi/g$ is a rank-one projector. Theorem~\ref{thm:certificate-sufficiency}\,(ii) states its
converse for exactly such block lists.

(iii) is the definition of the transpose; (i) and (ii) were proved from the
operator identities above and read no entry of $\Psi$ or of $M_x$. That
these are not declarable in general is seen on $S$ a qubit:
$\tau_S=\mathrm{diag}(\tfrac12,\tfrac12)$ gives $\Psi$ with every non-zero
entry equal to $\tfrac12$, while $\tau_S=\mathrm{diag}(\tfrac34,\tfrac14)$
gives $\Psi$ the entry $\sqrt{3/16}$, which is irrational.

(iv) In the proof of (SR2) the completion term contributed
$K(\idon{}-Q)K=0$, and the same holds of any operator supported on the
kernel of $\tau_S$; in (SR1) and (ii) $M_x$ does not occur; in (v) below it
contributes $\tau_A^{1/2}(\idon{}-P)\tau_A^{1/2}=0$.

(v) $\tau_A^{1/2}\tau_A^{-1/2}=P$, so
$\tau_A^{1/2}M_x\tau_A^{1/2}=P(\sigma^{\mathrm v}_x)^{\transpose}P
=\big(Q\,\sigma^{\mathrm v}_xQ\big)^{\transpose}
=(\sigma^{\mathrm v}_x)^{\transpose}$, using (c) and the support inclusion
above. The sum is $\tau_S^{\transpose}=\tau_A$.
\end{proof}
